\documentclass[11pt]{article}
\usepackage[margin=0.86in]{geometry}
\usepackage{amsmath,amssymb,amsthm,mathtools}
\usepackage{lmodern,microtype}
\usepackage[hidelinks]{hyperref}
\setlength{\emergencystretch}{2em}
\predisplaypenalty=10000
\title{A five-variable piecewise cubic function\\
without a polynomial max--min representation}
\author{}
\date{\small Draft for mathematical review}
\begin{document}
\maketitle
\vspace{-1em}

\section{The whole-space function}

Let \((a,b,c,d,e)\in\mathbb R^5\), and define
\[
x=a-6c+e,\quad y=4(d-b),\quad z=a+2c+e,
\quad u=2(e-a),\quad v=4(d+b),\quad R=2b+4d+2a^2.
\tag{1}
\]

The first five coordinates are an invertible linear change. Thus this is
the actual polynomial graph
\(R=(y+3v)/4+2((x+3z-2u)/8)^2\) over all five original coordinates.
Use
\[
C=\begin{pmatrix}
x&-v+R/2&-y\\
y&z+u&x\\
z&R/2&0\\
u&-v/2+R&-v/2\\
v&2z+u/2&u/2
\end{pmatrix},\qquad J=\operatorname{diag}(1,1,3,-1,-1).
\tag{2}
\]

Put \(Q=C^TJC\), and let \(D\) be the determinant of the first three
rows of \(C\). The full polynomial identities are
\[
Q=\begin{pmatrix}A_0&B_0&0\\B_0&E_0-A_0&F_0\\0&F_0&E_0\end{pmatrix},
\]
\[
A_0=16(ae-4bd+3c^2),\qquad
B_0=8(2cd-3be+a^3),
\]
\[
E_0=4ae-12ac-12ce+36c^2+12d^2-40bd+12b^2,
\]
\[
F_0=4ae+12ac-20ce-12c^2+12d^2-12b^2+8bd+8a^2b.
\tag{3}
\]

Define the four original polynomial labels
\[
q_1=-10B_0+2E_0+2F_0,\qquad
q_2=4A_0+2B_0-2E_0-2F_0,
\]
\[
q_3=-4A_0+2B_0+6E_0-10F_0,\qquad
q_4=-4A_0+2B_0+2E_0+2F_0.
\tag{4}
\]

Let \(h=\min(q_1,q_2,q_3,q_4)\), and set
\[
f=\begin{cases}h& h>0\text{ and }D>0,\\0&\text{otherwise.}\end{cases}
\tag{5}
\]

This is a continuous finite-piece polynomial function on all of
\(\mathbb R^5\). For completeness, the fixed simplex reconstruction
identity for matrices of the shape (3), with the labels (4), gives
\[
h>0\quad\Longrightarrow\quad Q\succeq (h/4)I.
\tag{6}
\]

One can verify (6) entirely by the following finite matrix identity. Set
\(s_i=4e_i-\mathbf1\) for \(i=0,1,2\), \(s_3=-\mathbf1\),
\(u_i=(e_i+\mathbf1)/4\) for \(i=0,1,2\), and \(u_3=0\).
For \(q_{ij}=-s_i^TQs_j\), the retained edges \(01,03,12,13\)
give \(q_1,q_2,q_3,q_4\), while
\[
q_{02}=2q_1+6q_2+5q_4,\qquad
q_{23}=q_1+2q_2+2q_4,
\]
\[
Q=\sum_{i<j}q_{ij}(u_i-u_j)(u_i-u_j)^T,
\qquad
\sum_{i<j}(u_i-u_j)(u_i-u_j)^T=(I+\mathbf1\mathbf1^T)/4.
\tag{7}
\]

If the first three rows of \(C\) annihilated a nonzero coefficient vector,
its Gram square would be minus the sum of the squares of the remaining
two row evaluations. Thus (6) implies \(D\ne0\). The explicit finite
closed semialgebraic cover is
\[
S_0=\{D\le0\}\cup\bigcup_j\{q_j\le0\},
\quad
S_i=\{D\ge0,\ q_j\ge0\ \forall j,\ q_i\le q_j\ \forall j\},
\tag{8}
\]

with labels \(0,q_1,\ldots,q_4\). The pieces cover all of space and
agree on overlaps; in particular, on \(D=0\) with nonnegative labels,
the minimum is zero. Finite closed pasting proves the asserted continuity.

\textbf{Theorem.} The function (5) is not a finite maximum of finite minima of
polynomials in \(\mathbb R[a,b,c,d,e]\), with arbitrary real coefficients
and unrestricted degrees. The invertible change (1) gives the same conclusion
in \(\mathbb R[x,y,z,u,v]\).

The proof below constructs, for every proposed finite collection of leaves,
two actual real points contradicting that representation. It does not
assume that a trace ideal is a separating ideal.

\section{Weighted identities and the base point}

Assign weights \((5,6,7,8,9)\) to \((a,b,c,d,e)\), respectively, and put
\[
S_\lambda(a,b,c,d,e)
 =(\lambda^5a,\lambda^6b,\lambda^7c,\lambda^8d,\lambda^9e).
\tag{9}
\]

The following three original polynomials are weighted homogeneous:
\[
A=ae-4bd+3c^2,\qquad
B=2cd-3be+a^3,\qquad
W=3d^2-4ce+a^2b,
\tag{10}
\]

with weights \(14,15,16\), respectively. Write
\[
g=b^2-ac,\qquad G=-a^2d+3abc-2b^3.
\tag{11}
\]

Their weights are \(12\) and \(18\). Directly from (3),
\[
\begin{split}
q_1(S_\lambda v)&=16\lambda^{14}A(v)-80\lambda^{15}B(v)
                                      +16\lambda^{16}W(v),\\
q_2(S_\lambda v)&=48\lambda^{14}A(v)+16\lambda^{15}B(v)
                                      -16\lambda^{16}W(v),\\
q_4(S_\lambda v)&=-48\lambda^{14}A(v)+16\lambda^{15}B(v)
                                      +16\lambda^{16}W(v),
\end{split}
\tag{12}
\]

and
\[
q_3(S_\lambda v)=192\lambda^{12}g(v)+O(\lambda^{14}),
\qquad
D(S_\lambda v)=8\lambda^{18}G(v)+O(\lambda^{20}).
\tag{13}
\]

The remainders in (13) are polynomial in \(\lambda\) and the coordinates
of \(v\), and are uniform for \(v\) in any fixed bounded set. For
the determinant assertion, a full expression is obtained by putting
\[
h_0=-ac-ce+6c^2+2d^2-4bd+2b^2,\quad
b_0=d(-a+4c+e)+2b(c-e),
\]
\[
D=8[(a+2c+e)b_0-(b+2d+a^2)h_0].
\tag{14}
\]

Its weight-18 part is exactly \(8G\), and every remaining monomial
has weight at least 20.

The monomial trace
\[
\tau(t)=(t^5,t^6,t^7,t^8,t^9)
\tag{15}
\]

satisfies \(Q(\tau(t))=0\) and \(D(\tau(t))=0\). We shall use its
single real base point
\[
p=\tau(1)=(1,1,1,1,1).
\tag{16}
\]

Weighted homogeneity, not an assumption of ordinary homogeneity of the
selection, is used in (9)--(13).

\section{Exact first and second variations}

All directions and corrections will have first coordinate zero, so their
affine points keep \(a=1\). On this affine hyperplane, \(A,B,W\)
are quadratic polynomials in the remaining four coordinates. For a direction
\(v=(0,v_b,v_c,v_d,v_e)\), write their homogeneous quadratic remainders as
\[
A_2(v)=-4v_bv_d+3v_c^2,
\quad B_2(v)=2v_cv_d-3v_bv_e,
\quad W_2(v)=3v_d^2-4v_cv_e.
\tag{17}
\]

Thus, for example,
\(A(p+\varepsilon v)=\varepsilon dA_p(v)+\varepsilon^2A_2(v)\)
exactly. Put
\[
F=W-B+A,\qquad F_2=W_2-B_2+A_2.
\tag{18}
\]

The full differential \(dF_p\) is zero. In particular,
\(F(p+\varepsilon v)=\varepsilon^2F_2(v)\) for all the stated directions.

Take the explicit rational vectors
\[
n=(0,2,3,3,2),\qquad j=(0,0,1,2,2).
\tag{19}
\]

They satisfy
\[
dA_p(n)=dB_p(n)=dW_p(n)=0,
\quad dA_p(j)=dB_p(j)=dW_p(j)=0,
\]
\[
dg_p(n)=1,\quad dg_p(j)=-1,\qquad
dG_p(n)=0,\quad dG_p(j)=1.
\tag{20}
\]

If \(F_2(v,w)\) denotes the symmetric polar form with
\(F_2(v,v)=F_2(v)\), then
\[
F_2(n)=0,\qquad F_2(n,j)=0,\qquad F_2(j)=3.
\tag{21}
\]

These identities follow immediately from (17):
\((A_2,B_2,W_2)(n)=(3,6,3)\),
\((A_2,B_2,W_2)(j)=(3,4,4)\), and their polar values on \((n,j)\)
are \((1,3,2)\).

Let \(e_d,e_e\) denote the corresponding coordinate unit vectors, and
use the three-dimensional correction space
\[
E=\operatorname{span}(e_d,e_e,j).
\tag{22}
\]

The first two derivative rows on this basis are
\[
\begin{pmatrix}-4&1&0\\2&-3&0\end{pmatrix}.
\tag{23}
\]

Their first two columns have determinant 10; the displayed vectors in (22)
are independent because only \(j\) has a nonzero \(c\)-coordinate.

\section{Actual corrected points with matched low-weight polynomials}

Fix any \(\delta>0\). Define
\[
s(\lambda)=3-3\lambda/8+\lambda^2,
\qquad r(\lambda)=\sqrt{s(\lambda)/3},
\qquad y_\pm(\lambda)=n\pm\delta r(\lambda)j.
\tag{24}
\]

Here \(s(\lambda)=(\lambda-3/16)^2+759/256>0\), so the square
root is smooth and \(r(0)=1\). For a correction \(z\in E\), impose
the three real equations
\[
\begin{split}
dA_p(z)+\varepsilon\bigl[A_2(y_\pm(\lambda)+z)
                                      -\lambda^2\delta^2\bigr]&=0,\\
dB_p(z)+\varepsilon\bigl[B_2(y_\pm(\lambda)+z)
                                      -(3\lambda/8)\delta^2\bigr]&=0,\\
F_2(y_\pm(\lambda)+z)-\delta^2s(\lambda)&=0.
\end{split}
\tag{25}
\]

At \((\varepsilon,\lambda,z)=(0,0,0)\) these equations hold. In the
basis (22), the derivative with respect to \(z\) has the form
\[
\begin{pmatrix}
-4&1&0\\2&-3&0\\ *&*&\pm6\delta
\end{pmatrix},
\tag{26}
\]

whose determinant is \(\pm60\delta\ne0\). The ordinary real implicit
function theorem therefore provides smooth actual solutions
\(z_\pm(\varepsilon,\lambda)\in E\) near \((0,0)\).
For every sufficiently small \(\lambda\), zero already solves (25) at
\(\varepsilon=0\), by (20), (21), (24); local uniqueness gives
\[
z_\pm(0,\lambda)=0.
\tag{27}
\]

Intersect the two parameter neighborhoods and choose a rectangle
\(|\varepsilon|<\varepsilon_0\), \(|\lambda|<\lambda_0\) contained
in that intersection. All subsequent choices of \(\varepsilon\) are
made inside this fixed rectangle. Hence the final step
\(\lambda\to0^+\), with \(\varepsilon\) already fixed, stays in
the domain of both actual solution functions.

Define actual real points
\[
v_\pm(\varepsilon,\lambda)
 =p+\varepsilon\bigl(y_\pm(\lambda)+z_\pm(\varepsilon,\lambda)\bigr).
\tag{28}
\]

The exact quadratic expansions, (25), and \(W=F+B-A\) imply
\[
A(v_\pm)=\lambda^2\varepsilon^2\delta^2,
\quad B(v_\pm)=(3\lambda/8)\varepsilon^2\delta^2,
\quad W(v_\pm)=3\varepsilon^2\delta^2.
\tag{29}
\]

There are no omitted Gram equations: only these three polynomials are being
prescribed, and the remaining label is checked separately below.

At \(\lambda=0\), smoothness and (20), (27) give
\[
g(v_\pm(\varepsilon,0))=\varepsilon(1\mp\delta)+O(\varepsilon^2),
\quad G(v_\pm(\varepsilon,0))=\pm\varepsilon\delta+O(\varepsilon^2).
\tag{30}
\]

Consequently, if \(0<\delta<1/2\) is fixed and \(\varepsilon>0\)
is sufficiently small, both values of \(g\) in (30) are positive, and
the two values of \(G\) have opposite signs. Fix such an \(\varepsilon\).
For sufficiently small positive \(\lambda\), put
\[
X_\pm=S_\lambda v_\pm(\varepsilon,\lambda).
\tag{31}
\]

Equations (12), (29) give exactly
\[
q_1(X_\pm)=34\lambda^{16}\varepsilon^2\delta^2,
\quad q_2(X_\pm)=q_4(X_\pm)=6\lambda^{16}\varepsilon^2\delta^2.
\tag{32}
\]

By (13), (30), and continuity, \(q_3(X_\pm)\) has a positive
weight-12 leading coefficient and is eventually larger than the values of
\(q_2,q_4\). The determinant signs are positive at \(X_+\) and
negative at \(X_-\). Thus
\[
f(X_+)=6\lambda^{16}\varepsilon^2\delta^2>0,
\qquad f(X_-)=0.
\tag{33}
\]

The choices are sequential: \(\delta\) is chosen first, then a fixed
positive \(\varepsilon\), and finally a sufficiently small positive
\(\lambda\). All points and equalities in (25)--(33) are over the reals.

\section{Every weighted-homogeneous test through weight 16}

Let \(\mathcal P_w\) be the vector space of weighted-homogeneous
polynomials of weight \(w\) for the weights (9). If \(H\in\mathcal P_w\),
then
\[
H(\tau(t))=t^w H(p).
\tag{34}
\]

For weights at most 16, all nonzero monomial spaces, apart from the constant,
are as follows:
\[
\begin{array}{c|l}
5,6,7,8,9&a,\ b,\ c,\ d,\ e\ \text{respectively}\\
10,11&a^2,\ ab\ \text{respectively}\\
12&ac,\ b^2\\
13&ad,\ bc\\
14&ae,\ bd,\ c^2\\
15&a^3,\ be,\ cd\\
16&a^2b,\ ce,\ d^2.
\end{array}
\tag{35}
\]

Weights 1--4 have no monomials. The list is exhaustive because any fourth
factor has weight at least 20, and the only possible triples through 16
are \(a^3\) and \(a^2b\).

When \(H(p)=0\), its coefficients in the appropriate row sum to zero.
The kernels and their derivative values on \(n\) are given by the following
explicit bases:
\[
\begin{array}{c|l|l}
w&\text{basis for }H(p)=0&\text{values of }dH_p(n)\\
12&b^2-ac&1\\
13&bc-ad&2\\
14&bd-ae,\ c^2-ae&3,\ 4\\
15&be-a^3,\ cd-a^3&4,\ 6\\
16&ce-a^2b,\ d^2-a^2b&3,\ 4.
\end{array}
\tag{36}
\]

For every other weight through 16, \(H(p)=0\) forces \(H=0\).
It follows from (36) that
\[
\{H\in\mathcal P_w:H(p)=0, dH_p(n)=0\}
=\begin{cases}
\mathbb R A&w=14,\\
\mathbb R B&w=15,\\
\mathbb R W&w=16,\\
\{0\}&0\le w\le16\text{ otherwise.}
\end{cases}
\tag{37}
\]

For example, \(A=3(c^2-ae)-4(bd-ae)\),
\(B=2(cd-a^3)-3(be-a^3)\), and
\(W=3(d^2-a^2b)-4(ce-a^2b)\). This calculation includes every real
weighted-homogeneous polynomial of the stated weights, not just the Gram
polynomials.

\section{A finite collection of arbitrary polynomial leaves}

Let \(\mathcal L\) be any finite collection of polynomials in
\(\mathbb R[a,b,c,d,e]\), with no restriction on coefficients or degrees.
We will choose the parameters in (31) so that
\[
P(X_+)\ge f(X_+)\quad\Longrightarrow\quad P(X_-)>0
\qquad(P\in\mathcal L).
\tag{38}
\]

Decompose each leaf into its finitely many weighted-homogeneous components,
\(P=\sum_w P_w\). Call a component through weight 16 exceptional if
it belongs to the corresponding space on the right of (37), including zero.
There are finitely many nonexceptional components through weight 16.

For such a component \(H\), either \(H(p)\ne0\), or
\(H(p)=0\) and \(dH_p(n)\ne0\), by (37). In the latter case,
choose \(\delta>0\) small enough that
\[
dH_p(n+\delta j),\qquad dH_p(n-\delta j)
\]

are both nonzero and have the sign of \(dH_p(n)\). One fixed
\(0<\delta<1/2\) works for all the finitely many components.
For \(\varepsilon>0\) sufficiently small, their values at
\(v_\pm(\varepsilon,0)\) consequently have the same nonzero sign:
if \(H(p)\ne0\), their leading term is that constant; otherwise it
is \(\varepsilon dH_p(n\pm\delta j)\). Choose one such
\(\varepsilon\) also satisfying (30).

The same choice can ensure the following slightly stronger statement when
needed. Adding any of the finitely many constants
\(c\varepsilon^2\delta^2\) arising below to these paired values does
not change their signs. Indeed the nonzero leading term is either a fixed
constant or is linear in \(\varepsilon\); the added term is quadratic.
The constants, components, and \(\delta\) are all fixed before choosing
this final \(\varepsilon\).

Now analyze one leaf \(P\). If it has a nonexceptional component through
weight 16, let \(w\) be the smallest such weight and write \(H=P_w\).
Every earlier nonzero component is a multiple of \(A\) or \(B\), and
by (29) its evaluation after (9) has effective order \(\lambda^{16}\).

If \(w<16\), it follows that
\[
\lambda^{-w}P(X_\pm)\longrightarrow H(v_\pm(\varepsilon,0))
\qquad(\lambda\to0^+).
\tag{39}
\]

The two limits have the same nonzero sign. If they are negative, the
hypothesis in (38) is eventually impossible; if they are positive,
the conclusion in (38) holds.

If \(w=16\), write the earlier components as \(\alpha A\) at
weight 14 and \(\beta B\) at weight 15. Formula (29) gives instead
\[
\lambda^{-16}P(X_\pm)\longrightarrow
H(v_\pm(\varepsilon,0))
 +(\alpha+3\beta/8)\varepsilon^2\delta^2.
\tag{40}
\]

The preceding strengthened choice of \(\varepsilon\) makes these limits
same-signed and nonzero. The same reasoning proves (38).

Finally, suppose all components through weight 16 are exceptional. Write
them as \(\alpha A,\beta B,\gamma W\) in their respective weights.
Then
\[
\lambda^{-16}P(X_\pm)\longrightarrow
c_P\varepsilon^2\delta^2,
\qquad c_P=\alpha+3\beta/8+3\gamma.
\tag{41}
\]

If \(c_P>0\), then \(P(X_-)>0\) for sufficiently small positive
\(\lambda\). If \(c_P\le0\), (33) and (41) give
\(P(X_+)<f(X_+)\) eventually, including the case \(c_P=0\).
Thus (38) also holds in this last case.

All components of weight greater than 16 contribute
\(o(\lambda^{16})\): the vectors \(v_\pm(\varepsilon,\lambda)\)
stay bounded and converge as \(\lambda\to0\), and there are only
finitely many components of each fixed leaf. No bound on their weights is
assumed. Since \(\mathcal L\) is finite, one sufficiently small
positive \(\lambda\) makes (33) and (38) hold simultaneously for every
leaf. This completes the unrestricted-degree finite-leaf assertion.

\section{Contradiction to a polynomial lattice representation}

Suppose that on all of \(\mathbb R^5\),
\[
f=\max_{1\le i\le m}\min_{1\le j\le n_i}P_{ij}
\]

for finitely many original real polynomials, with \(m,n_i\ge1\).
Apply Section 6 to that finite set of leaves. At \(X_+\), a row attains
the outer maximum \(f(X_+)>0\); every leaf in that row is at least
\(f(X_+)\). By (38), every one of those finitely many leaves is
strictly positive at \(X_-\). Their finite minimum is strictly positive,
so the represented value at \(X_-\) is strictly positive. This contradicts
the actual value \(f(X_-)=0\) in (33).

This proves the theorem for arbitrary real coefficients and arbitrary finite
degrees. Any finite expression using minima and maxima can be converted to
the displayed form by distributivity, so allowing a different finite lattice
syntax does not change the conclusion.

The construction is five-dimensional because its free input is exactly
\((a,b,c,d,e)\), equivalently the five coordinates \((x,y,z,u,v)\) in
(1). It does not prove that dimension five is minimal or settle any universal
positive statement in dimensions three or four.

\end{document}
